\section{Pose as a Stable Virtual Camera Signal}
\label{sec:pose}

Tracking answers \emph{which} surfer is present in each frame, but a useful focused
view also needs to decide where to place the virtual camera and how tightly to crop.
A detector bounding box is a poor control signal for this purpose.  A windsurf rig
is articulated, strongly anisotropic, and continually rotates with respect to the
camera.  Its axis-aligned box can therefore change width, height, and centre even
when the rider has moved little and remains at approximately the same distance.
Using the box centre and height directly makes the crop drift across the subject and
causes visible zoom pumping.

\subsection{Two semantic points instead of a generic human pose}

The production detector is a YOLO pose model supervised for one class and two
domain-specific keypoints: the boom--mast junction and the mast tip.  These points
were chosen for the rendering geometry rather than to describe the complete human
body.  The boom--mast junction gives a repeatable reference near the rider and sail,
while the distance from that junction to the mast tip provides a rig-relative scale cue
that is substantially less sensitive than an axis-aligned box height to in-plane
rotation.

The labels build on the manually reviewed detection boxes.  A dedicated annotation
tool presents an enlarged crop for each box, permits either keypoint to be marked as
not visible, and writes two-keypoint YOLO pose labels.  The checked-in pose project
contains 1,123 manually labelled training images and 56 validation images.  The
workflow also supports conservative pseudo-labelling: a prediction is accepted only
after matching it to an existing box and applying detection, overlap, keypoint
confidence, and spatial plausibility gates; pseudo-labels remain separate from manual
labels and are included in training only when requested.  Finally, a hard-window
miner runs the complete local pipeline and ranks short intervals by keypoint jitter
relative to the rendered anchor.  This directs subsequent review towards failures
that are visually disruptive in the final crop, rather than towards another uniform
sample of easy frames.

\subsection{From keypoints to crop centre}

For a detection with bounding box top $y_1$, height $h$, and centre
$(c_x,c_y)$, the fallback anchor is
\begin{equation}
  \mathbf{a}_{\mathrm{box}} = (c_x,\; y_1 + 0.6h).
  \label{eq:pose-box-anchor}
\end{equation}
Placing it below the geometric box centre is a deliberate bias towards the rider and
board rather than the empty upper part of a tall sail box.  When the boom--mast
junction $\mathbf{b}=(b_x,b_y)$ is visible, the implementation first constructs a
low-noise proxy for the mast tip, $\tilde{\mathbf{t}}=(b_x,y_1)$, and places the
semantic anchor at
\begin{equation}
  \mathbf{a}_{\mathrm{pose}}
    = 0.15\tilde{\mathbf{t}} + 0.85\mathbf{b}.
  \label{eq:pose-semantic-anchor}
\end{equation}
Thus the crop is anchored slightly above the boom junction.  The measured mast-tip
coordinate is intentionally not used for position: its larger lever arm would turn
localization noise into unnecessary crop motion.  It remains valuable for scale.

\begin{figure}[t]
  \centering
  \includegraphics[width=0.58\linewidth]{pose-anchor-geometry.png}
  \caption{Pose-derived virtual-camera geometry on a representative labelled
  frame.  The detector box and its centre (orange) describe image extent.  The
  boom--mast junction and mast tip (cyan) instead provide a rider-relative anchor
  (magenta) and a rig-relative scale cue.  Production rendering smooths these
  signals over time and falls back to box geometry when keypoints are unavailable.}
  \label{fig:pose-anchor-geometry}
\end{figure}

Small detections do not contain enough pixels for equally reliable keypoint
localization.  The final unsmoothed anchor therefore blends
$\mathbf{a}_{\mathrm{pose}}$ with $\mathbf{a}_{\mathrm{box}}$ according to the
shorter box dimension.  The pose weight rises linearly from zero at 45 pixels to one
at 55 pixels.  A keypoint is considered visible at confidence 0.15 or greater; below
that threshold, the box fallback is used.

Missing observations are handled as transitions, not abrupt mode switches.  Gaps of
up to five frames between visible boom keypoints are interpolated directly between
semantic anchors.  Across longer gaps, the anchor eases from the last semantic point
to the per-frame box fallback and eases back when a reliable keypoint returns.
Leading and trailing missing spans receive analogous fallback or short hold
behaviour.  A final five-frame robust weighted mean suppresses isolated jumps while
giving the current frame the largest weight; Huber reweighting limits the influence
of spatial outliers.  This processing operates on the dense tracks produced by the
tracking post-processor, so every frame in a track receives an anchor even when its
detection was interpolated.

\subsection{A scale tied to the rig}

When both keypoints are visible, the scale signal is the Euclidean length
$\ell=\lVert\mathbf{b}-\mathbf{t}\rVert_2$ of the observed boom-to-tip mast
segment.  Short missing spans are interpolated and the last valid value is held for
up to five frames; otherwise the vertical distance from the current anchor to the
box top supplies a fallback.  A centred moving mean with radius 15 frames smooths
this length signal.  The requested crop height in pixels is then
\begin{equation}
  C = \max\!\left(\frac{\bar{\ell}}{0.75},
                  \frac{h}{0.9}, 1\right),
  \qquad s = \operatorname{clip}\!\left(\frac{C}{H},10^{-6},1\right),
  \label{eq:pose-scale}
\end{equation}
where $H$ is the source-video height.  The first term aims to make the semantic mast
segment occupy 75\% of the output height; the second prevents the detected box from
being cropped more tightly than a 90\% fill.  Taking the maximum preserves context
when either signal calls for a wider view.

The resulting normalized anchor and scale are stored with every renderable
detection and consumed directly by the focused player.  They stabilize the
\emph{composition} of an individual rider: which semantic point stays near the
centre and how large the rig appears.  This is distinct from global camera-motion
compensation.  The latter estimates frame-to-frame background transforms (while
masking detected surfers), supports motion-based association, and supplies an
optional stabilized overview.  It does not replace the pose-derived virtual camera;
a globally stabilized frame can still contain a rotating rig whose bounding box
centre and dimensions fluctuate.
